(1) 基模高斯光束的Rayleigh距离为 $z_R=\pi w_0^2/\lambda_0$，是束腰处横向衍射尺寸与轴向传播变化的特征长度。对应的波数 $k=2\pi/\lambda_0$ 可写 $z_R=kw_0^2/2$，其量纲确为长度。
(2) 近轴Helmholtz方程的高斯包络可用复光束参数 $q(z)=z+iz_R$ 描述，光束强度半径满足 $w(z)=w_0\sqrt{1+(z/z_R)^2}$。在束腰平面 $z=0$ 有最小半径 $w_0$，左右传播方向的光束半径对 $z$ 为偶函数。
(3) 当 $|z|\gg z_R$，$w(z)\simeq w_0|z|/z_R$，所以远场发散半角 $\theta_{\rm div}\simeq w_0/z_R=\lambda_0/(\pi w_0)$。在Rayleigh位置 $z=z_R$，$w(z_R)=\sqrt2\,w_0$，强度横向截面积随之约增一倍。
(4) 当束腰 $w_0$ 减小时，$z_R\propto w_0^2$ 显著缩短，而发散角 $\theta_{\rm div}\propto1/w_0$ 增大。若 $w_0$ 与光波长同量级，发散角不再足够小，简单近轴高斯公式失效；几何光线不含这类由有限横向孔径产生的衍射发散。
